\newif\ifglobal

%%%%%%%%%%%%%%%%%%%%%%%
%%% Compile to view %%%
%%%%%%%%%%%%%%%%%%%%%%%

%%%%%%%%%%%%%%%%%%%%%%%%%%%%%%%%%%%%%%%%%%%%%%%%%%%
%%% Readme:                                     %%%
%%% https://www.overleaf.com/read/bwdjvfjksvjy/ %%%
%%%%%%%%%%%%%%%%%%%%%%%%%%%%%%%%%%%%%%%%%%%%%%%%%%%

% >>> M?: marks a module setting number or important notice

% >>> M4: Set {./relative-path} to external files for this module <<<
\def \modulefiles {./61-TCHSEN}
% To add an external file, use:
% (Replace '---' with actual filename)
% '\input{\modulefiles/tables/---.tex}' for tables
% '\includegraphics{\modulefiles/figures/---}' for figures
% '\subfile{\modulefiles/---}' for register files

% >>> M1: This module is in global project? <<<
% 'YES' - uncomment; 'NO' - comment out
\globaltrue

\ifglobal
    \documentclass[main\_\_EN.tex]{subfiles}
   
\else
    \documentclass[openany, 10pt]{book}

    % >>> M2: In ./00-shared/preamble-trm-module, also find
    %         settings for visibility of todo notes, watermark <<<
    \usepackage{preamble-trm-module}

    % >>> M3: Temporary labels in English,
    %         you can add your own labels there <<<
    \myexternaldocument{./00-shared/temp-labels-en}

\fi %\ifglobal

\setcounter{secnumdepth}{4}

% >>> M5: If this module needs any updates in
% './00-shared/preamble-trm-module.sty'
% (include additional package, etc.), contact your module PM
% or create the file './00-shared/changelog.tex' make the
% necessary updates yourself and document them in this file <<<

\raggedright


\begin{document}


% Sets language into which pre-defined bits of text are translated
\selectlanguage{english}

% Do not delete this block
\ifglobal
% Add the following front matter if
% this module is outside of global project
\else
    \InputIfFileExists{readme}{}{}
    {\let\clearpage\relax \listoftodos}
    {\let\clearpage\relax \listofdonetodos}
    \tableofcontents
    %\listoftables
    %\listoffigures
\fi


%%%%%%%%%%%%%%%%%%%%%%%%%
%%% TRM Module Begins %%%
%%%%%%%%%%%%%%%%%%%%%%%%%

\hypertarget{tchsen}{}
\chapter{Touch Sensor (TOUCH)}
\label{mod:tchsen}

\chipname{} has 14 capacitive touch sensor units, each of which can be connected to an external touch panel via a touch pin (GPIO pin) to constitute a touch sensor system. The \chipname{} touch sensor system is mainly used in human-computer interactions to detect finger touch or proximity. It also features moisture resistance and waterproof design.

\section{Terminology}
\label{sec:tchsen-term}

To better illustrate the functions of capacitive touch sensors, the following terms are used in this section.

\begin{longtable}[c]{ R{7cm} p{9cm} }
\textbf{Touch Sensor}  & Touch-related internal sensing circuitry. \\
\textbf{Touch Pin}    & Pins with touch sensing feature. \\ 
\textbf{Touch Panel}  & The external device connected to touch sensor via channel (touch pin) to detect finger touch. \\
\textbf{Touch Sensor System}  & Capacitive touch sensing system, consisting of touch sensor, touch pin, traces, and touch panel. \\
\end{longtable}

\vspace{-1em}
\begin{tiplisting}
    In subsequent description, ``touch panel is sampled/scanned/measured" indicates the same action to touch pin, i.e. ``touch pin is sampled/scanned/measured".
\end{tiplisting}
\vspace{-2em}


%%%%%%%%%%%%%%%%%%%%
%%% Feature List %%%
%%%%%%%%%%%%%%%%%%%%

\section{Feature List}
\label{sec:tchsen-features}

The \chipname{} touch sensor has the following features:

\begin{itemize}
    \item Detection of 14 capacitive touch pins
    \item Sampling triggered by software or dedicated hardware timer
    \item Two sampling methods:
        \begin{itemize}
            \item Pulses from the touch pins used as clock signals to count the sampling period
            \item Pulses from the touch pins used as digital signals; sample the rising edge of the digital signal with the system clock to count the sampling period
        \end{itemize}
    \item Scan mode, supporting sequential sampling of multiple touch pins by configuring the Touch FSM.
    \item Timeout mechanism to monitor channel abnormality
    \item Frequency hopping to increase the anti-interference of detection
    \item Proximity sensing mode with up to three configurable channels
    \item Configuration of individual touch sensors to operate normally in sleep mode
    \item Wake-up by touch sensor
    \item Moisture resistance
    \item Waterproof design
\end{itemize}


%%%%%%%%%%%%%%%%%%%%%%%%%%%%%%
%%% Architectural Overview %%%
%%%%%%%%%%%%%%%%%%%%%%%%%%%%%%

\section{Architectural Overview}
\label{sec:tchsen-arch-overview}

\subsection{Touch Panel}

A touch panel consists of the following components:

\begin{itemize}
    \item Electrode: Changes in capacitance when touched by a finger.
    \item Substrate: The base material on which the protective cover, electrode, and the electrode's connector to the channel are built.
    \item Protective cover: To physically separate the other components from the external environment.
\end{itemize}

When developing applications using the touch sensing feature, users should decide on the placement, the material, or the arrangement of the touch panels during mechanical design. For more information about the design guidelines, please refer to \href{https://github.com/espressif/esp-iot-solution/blob/release/v1.1/documents/touch_pad_solution/touch_sensor_design_en.md}{Touch Sensor Application Note}. Touch panel can be connected to \chipname{} touch sensor via touch pin (channel). See Figure \ref{figure:touch sensor}.

\begin{figure}[H]
    \centering
    \fbox{\includegraphics[width=0.9\textwidth]{\modulefiles/figures/touch-sensor-structure.pdf}}
    \caption{Touch Sensor System}
    \label{figure:touch sensor}
\end{figure}

When users touch the protective cover, the capacitance of the electrode increases. If the electrode is connected to one of the \chipname{} touch sensors via a channel, the touch sensor will be able to detect the electrode's change in capacitance. If the change of capacitance exceeds the configurable threshold, the touch sensor will trigger an interrupt.

\subsection{Capacitive Touch Pin}

\chipname{} provides 14 capacitive touch pins labeled as T1 \verb+~+ T14, each connected to the chip's pin to monitor finger touch from the external environment. The mapping between the touch sensors and the chip pins is shown in Table \ref{tab:touch-pads}.

\begin{longtable}[c]{|p{4.0cm}|p{6.0cm}| }
    \caption{\chipname{} Capacitive Touch Pins}
    \label{tab:touch-pads}
    \\\hline
    \rowcolor{lightgray}
    \textbf{Capacitive Touch Pin} & \textbf{Connected GPIO Pin} \\ 
    \hline
    \endfirsthead
    \hline
    \rowcolor{lightgray}
    \textbf{Capacitive Touch Pin} & \textbf{Connected GPIO Pin} \\  
    \hline
    \endhead
    \endfoot
    \endlastfoot
    T1  &  GPIO2 \\\hline
    T2  &  GPIO3 \\\hline
    T3  &  GPIO4  \\\hline
    T4  &  GPIO5  \\\hline
    T5  &  GPIO6  \\\hline
    T6  &  GPIO7  \\\hline
    T7  &  GPIO8 \\\hline
    T8  &  GPIO9 \\\hline
    T9  &  GPIO10 \\\hline
    T10  &  GPIO11 \\\hline
    T11 &  GPIO12 \\\hline
    T12 &  GPIO13 \\\hline
    T13 &  GPIO14 \\\hline
    T14 &  GPIO15 \\\hline
\end{longtable}

\subsection{Touch Sensor}

The touch sensor charges and discharges the touch pin with a fixed current source, and each charge and discharge generates a pulse signal, the TOUCH\_OUT signal. The touch sensor detects the capacitance change of the touch pin. If the touch pin is touched or approached by a finger, the touch pin's capacitance increases and the charging and discharging time is extended. By measuring the time required for charging and discharging the touch pin a fixed number of times, it can be inferred whether the touch pin has been touched or not.

Figure \ref{fig:touch-sensor-internal-structure} illustrates the internal structure of the touch sensors. Each of the touch sensors has a set of input and output signals, some of which are connected to the Touch FSM (see Section \ref{sec:touch-fsm}). Table \ref{tab:tchsen-output-signal} shows the output signals of the touch sensor and their functions.

\begin{figure}[H]
    \centering
    \includegraphics[width=1.0\textwidth]{\modulefiles/figures/tchsen-internal-structure}
    \caption{Touch Sensor Internal Structure}
    \label{fig:touch-sensor-internal-structure}
\end{figure}

\begin{longtable}[c]{|p{3cm}|p{2cm}|p{11cm}| }
    \caption{Capacitive Touch Sensor Output Signal}
    \label{tab:tchsen-output-signal}
    \\\hline
    \rowcolor{lightgray}
    \textbf{Signal} & \textbf{Source}& \textbf{Function}\\ 
    \hline
    \endfirsthead
    \hline
    \rowcolor{lightgray}
    \textbf{Signal} & \textbf{Source}& \textbf{Function}\\ 
    \hline
    \endhead
    \endfoot
    \endlastfoot
    TOUCH\_XPD & Touch FSM & Controls the power-up and power-down of the touch sensor \\\hline
    TOUCH\_START & Touch FSM & Controls the touch sensor to initiate a measurement. The measurement ends when the Touch FSM pulls down the TOUCH\_START signal \\\hline
    DBIAS & Touch FSM & The FSM selects different DBIAS based on the frequency mode to adjust the voltage of the touch sensor's internal LDO, which affects the charging and discharging speed of the touch sensor \\\hline
    DRV\_HS & Touch FSM & The FSM selects different DRV\_HS based on the frequency mode to adjust the high-speed driving capability of the touch sensor \\\hline
    DRV\_LS & Touch FSM & The FSM selects different DRV\_LS based on the frequency mode to adjust the low-speed driving capability of the touch sensor \\\hline
    DRES\_LPF & Touch FSM & The FSM selects different DRES\_LPF based on the frequency mode to adjust the RC low-pass filter, which can be configured as 0 k$\Omega$, 3 k$\Omega$, 4.5 k$\Omega$, or 7.5 k$\Omega$ \\\hline
    DCAP\_LPF & Touch FSM & The FSM selects different DCAP\_LPF based on the frequency mode to adjust the RC low-pass filter with an adjustment range of 0-2.54 pF and a step of 20 fF \\\hline
    DCAP\_DRV & Touch FSM & Used to adjust the internal capacitance connected to touch channels during measurement, usually set to 0 \\\hline
    EN\_CAL & Touch FSM & Used to disconnect from the touch pins and to test the capacitance in the touch sensor (the touch panel needs to be disabled) \\\hline
    BUF\_DRV & Touch FSM & Adjusts the driving capability of TOUCH BUFF\\\hline
    EN\_BUF & Touch FSM & Enables TOUCH BUFF, a software-configurable touch panel, to mitigate the effects of water droplets on the sensor for moisture tolerance \\\hline
    TOUCH\_OUT & Touch sensor & The Touch FSM counts the pulses in the TOUCH\_OUT signal and measures the time it takes to generate N pulses \\\hline
\end{longtable}


%%%%%%%%%%%%%%%%%%%%%%%%%%%%%%
%%% Functional Description %%%
%%%%%%%%%%%%%%%%%%%%%%%%%%%%%%

\section{Functional Description}
\label{sec:tchsen-func-descr}

\subsection{Touch FSM}
\label{sec:touch-fsm}

The Touch FSM performs a measurement by selecting a touch sensor and controlling the necessary signals to and from a touch sensor. Figure \ref{fig:touch-fsm} shows the internal structure and operating mechanism of the Touch FSM. The Touch FSM is clocked by LP\_DYN\_FAST\_CLK. For more information on LP\_DYN\_FAST\_CLK, see Chapter \ref{mod:resclk} \textit{\nameref{mod:resclk}}.

\begin{figure}[H]
    \centering
    \fbox{\includegraphics[width=1.0\textwidth]{\modulefiles/figures/touch_fsm.pdf}}
    \caption{Touch FSM Structure and Mechanism}
    \label{fig:touch-fsm}
\end{figure}

The following introduces the modules in the Touch FSM.

\begin{itemize}
    \item SCAN\_CTRL: Select the touch pin to measure, set the parameters such as hopping frequency, and controls the start and end of measurement for each touch sensor in scan mode.
    \item WORK\_UNIT: Sample the selected touch sensors during a measurement.
    \item Filter: If enabled, each touch sensor can filter a series of measurements via an infinite impulse response (IIR) filter. The filtered value will be returned as the sampled value. See Section \ref{para:sampled-value} for details.
    \item SCAN\_CUR: Cache the parameters of different touch sensors at different frequency modes during the measurement to be used for the next measurement.
\end{itemize}

\subsection{Sampled Signal Preprocessing}
\label{sec:sample-mode}

The core mechanism of the touch sensor is to reflect whether the touch panel is touched or not based on the change of the output pulse wave, i.e., the TOUCH\_OUT signal, through detecting the change of capacitance on the touch panel. Therefore, the sampling of the TOUCH\_OUT signal is critical. Figure \ref{fig:touch-out-sample} shows the preprocessing of the TOUCH\_OUT signal.

\begin{figure}[H]
    \centering
    \fbox{\includegraphics[width=1.0\textwidth]{\modulefiles/figures/touch_data_div.pdf}}
    \caption{TOUCH\_OUT Signal Preprocessing}\label{fig:touch-out-sample}
\end{figure}

The frequency of TOUCH\_OUT might be much higher than the sampling frequency. Therefore, in order to improve the accuracy of the signal sampling, users can divide the TOUCH\_OUT signal by configuring \hyperref[fielddesc:LPANADIVNUM0]{LP\_ANA\_DIV\_NUM\regindex{n}}. \regindex{n} in this field represents different frequency modes in frequency hopping (refer to Section \ref{sec:touch-freq-hop}) with the value 0-2. The lowest bit configures whether to enable or disable frequency division. Set the lowest bit to 0 to disable frequency division, and set it to 1 to enable division. The highest two bits indicate the frequency division coefficient with a value of 0/1/2/3, corresponding to the 0/2/4/6 divisions.

If the frequency is still higher than twice the sampling frequency after frequency division, set \hyperref[fielddesc:LPANATOUCHOUTSEL]{LP\_ANA\_TOUCH\_OUT\_SEL} to 1 to directly use the divided TOUCH\_OUT signal as a clock. Set \hyperref[fielddesc:LPANATOUCHOUTGATE]{LP\_ANA\_TOUCH\_OUT\_GATE} to configure whether to enable this clock.

\subsubsection{Measurement Process}
\label{para:measurement-process}

Depending on whether frequency hopping is enabled, a measurement of a touch pin can be divided into the following two processes:

\begin{itemize}
    \item Frequency hopping not enabled:
        \begin{enumerate}
            \item The Touch FSM selects the touch sensor to be measured. The relevant signals are routed to that touch sensor.
            \item The Touch FSM drives the START signal to the touch sensor to initiate the measurement. Internally, the Touch FSM starts a touch counter to time the duration of the measurement.
            \item A pulse counter in the Touch FSM is incremented on each pulse received from the touch sensor's TOUCH\_OUT signal.
            \item When the pulse counter reaches the count threshold set in \hyperref[fielddesc:LPANATOUCHMEASNUM0]{LP\_ANA\_TOUCH\_MEAS\_NUM0}, the measurement is complete. The START signal is de-asserted, and the touch counter is stopped. The value of the stopped touch counter x LP\_DYN\_SLOW\_CLK cycles indicates the time taken to charge and discharge the touch pin for \hyperref[fielddesc:LPANATOUCHMEASNUM0]{LP\_ANA\_TOUCH\_MEAS\_NUM0} cycles.
        \end{enumerate}
    \item Frequency hopping enabled:
        \begin{enumerate}
            \item The Touch FSM selects the touch sensor to be measured. The relevant signals are routed to that touch sensor.
            \item Set \hyperref[fielddesc:LPANAFREQSCANEN]{LP\_ANA\_FREQ\_SCAN\_EN} to 1 to enable frequency hopping, configure \hyperref[fielddesc:LPANAFREQSCANCNTLIMIT]{LP\_ANA\_FREQ\_SCAN\_CNT\_LIMIT} to set the number of frequency modes supported for frequency hopping (up to three), and configure the parameters required for each frequency mode and the number of pulses for sampling the TOUCH\_OUT signal.
            \item The Touch FSM drives the START signal to the touch sensor to initiate the measurement. Internally, the Touch FSM starts a touch counter to time the duration of the measurement.
            \item A pulse counter in the Touch FSM is incremented on each pulse received from the touch sensor's TOUCH\_OUT signal.
            \item When the pulse counter reaches the count threshold set in \hyperref[fielddesc:LPANATOUCHMEASNUM0]{LP\_ANA\_TOUCH\_MEAS\_NUM0} of the corresponding frequency mode, the measurement is complete. The value of the stopped touch counter indicates the time taken to charge and discharge the touch pin \hyperref[fielddesc:LPANATOUCHMEASNUM0]{LP\_ANA\_TOUCH\_MEAS\_NUM0} cycles with the current frequency.
            \item The sampling frequency counter is incremented by 1 and the touch counter is cleared. The next frequency mode will be sampled after a small delay. When the samplings are finished for all frequency modes, the START signal is de-asserted, and the touch counter is cleared. 
        \end{enumerate}
\end{itemize}

The sampled value (measured value) is the value of the touch counter indicated as touch\_raw\_data. The value of touch\_raw\_data can be read from \hyperref[fielddesc:RTCTOUCHPAD0DATA]{RTC\_TOUCH\_PAD\regindex{n}\_DATA}. Note that the value returned by \hyperref[fielddesc:RTCTOUCHPAD0DATA]{RTC\_TOUCH\_PAD\regindex{n}\_DATA} may also be the filtered touch\_raw\_data, i.e., touch\_smooth\_data and benchmark. Configure \hyperref[fielddesc:LPANATOUCHDATASEL]{LP\_ANA\_TOUCH\_DATA\_SEL} and \hyperref[fielddesc:LPANATOUCHFREQSEL]{LP\_ANA\_TOUCH\_FREQ \_SEL} to select the type of the return value. For more information on the sampled value types and how to judge touch actions based on the data, refer to Section \ref{para:sampled-value}.

\begin{tiplisting}
Note that if the pulse does not reach the set threshold for a long time while the touch counter has reached the timeout threshold set in \hyperref[fielddesc:LPANATOUCHTIMEOUTNUM]{LP\_ANA\_TOUCH\_TIMEOUT\_NUM} with \hyperref[fielddesc:LPANATOUCHTIMEOUTEN]{LP\_ANA\_TOUCH\_TIMEOUT\_EN} set as 1, the \hyperref[fielddesc:RTCTOUCHTIMEOUTINTRAW]{ RTC\_TOUCH\_\\TIMEOUT\_INT\_RAW} timeout interrupt will be triggered, indicating a circuit abnormality.
\end{tiplisting}

\subsubsection{Trigger Source of Measurement}

The Touch FSM initiates a measurement by sending a START signal. The START signal can either be triggered by software, or by a dedicated hardware timer known as the ``touch timer". The use of a touch timer allows periodic measurements to be taken without software intervention.

The touch timer is clocked by LP\_DYN\_SLOW\_CLK and should be configured with a period in number of LP\_DYN\_SLOW\_CLK cycles. The START signal is generated when the touch timer expires. The touch timer is reset when the measurement completes and restarts counting until the next expiration.

\begin{itemize}
    \item To trigger the START signal by software:
    \begin{itemize}
        \item Set \hyperref[fielddesc:LPANATOUCHSTARTFORCE]{LP\_ANA\_TOUCH\_START\_FORCE} to 1 to trigger the START signal by software.
        \item Software sets \hyperref[fielddesc:LPANATOUCHSTARTEN]{LP\_ANA\_TOUCH\_START\_EN} to 1 and generates the START signal to initiate a measurement.
    \end{itemize}
    \item To trigger the START signal by the touch timer:
    \begin{itemize}
        \item Clear \hyperref[fielddesc:LPANATOUCHSTARTFORCE]{LP\_ANA\_TOUCH\_START\_FORCE}.
        \item Configure the wait time before the timer sends the START signal (unit: LP\_DYN\_SLOW\_CLK cycle) via \hyperref[fielddesc:PMUTOUCHWAITCYCLES]{PMU\_TOUCH\_WAIT\_CYCLES}.
        \item Configure the wait time after the timer receives the DONE signal upon the completion of the touch sensor sampling (unit: LP\_DYN\_SLOW\_CLK cycle) via \hyperref[fielddesc:PMUTOUCHSLEEPCYCLES]{PMU\_TOUCH\_SLEEP\_CYCLES}.
        \item Set \hyperref[fielddesc:PMUTOUCHSLEEPTIMEREN]{PMU\_TOUCH\_SLEEP\_TIMER\_EN} to 1 to enable the touch timer.
    \end{itemize}
    \item To force exit the current measurement cycle of the touch timer by software, there are two approaches:
    \begin{enumerate}
        \item Generate the end signal of the touch sensor to exit the current measurement cycle of the touch timer.
        \begin{itemize}
            \item Set \hyperref[fielddesc:LPANATOUCHDONEFORCE]{LP\_ANA\_TOUCH\_DONE\_FORCE} to 1 to select the DONE signal to be triggered by software.
            \item Software sets \hyperref[fielddesc:LPANATOUCHDONEEN]{LP\_ANA\_TOUCH\_DONE\_EN} to 1 to generate the DONE signal to end a measurement.
        \end{itemize}
        \item Use the software to directly force the touch timer to exit the current measurement cycle.
        \begin{itemize}
            \item Set \hyperref[fielddesc:PMUTOUCHFORCEDONE]{PMU\_TOUCH\_FORCE\_DONE} to generate the DONE signal and force the touch timer to exit the current measurement cycle.
        \end{itemize}
    \end{enumerate}

\begin{tiplisting}
The touch timer waits for a period of time before generating the START signal and after receiving the DONE signal, so the interval between two consecutive measurements is (\hyperref[fielddesc:PMUTOUCHWAITCYCLES]{PMU\_TOUCH\_WAIT\_CYCLES} + \hyperref[fielddesc:PMUTOUCHSLEEPCYCLES]{PMU\_TOUCH\_SLEEP\_\\CYCLES})，and the timing clock is LP\_DYN\_SLOW\_CLK.
\end{tiplisting}
\end{itemize}

\subsubsection{Scan Mode}

Scan mode involves the Touch FSM taking measurements of multiple touch sensors in sequential order. On every START signal, a new touch sensor is selected for measurement, thus allowing multiple touch pins to be monitored. The scan process is illustrated in Figure \ref{fig:touch-scan}.

\begin{figure}[H]
    \centering
    \includegraphics[width=0.8\textwidth]{\modulefiles/figures/touch-scan}
    \caption{Timing Diagram of Scan Mode}
    \label{fig:touch-scan}
\end{figure}

To enable the scan mode, specify the bit map of the enabled touch pins in \hyperref[fielddesc:LPANATOUCHSCANPADMAP]{LP\_ANA\_TOUCH\_SCAN\_PAD\_MAP}. The Touch FSM will select one touch pin from the bit map of the enabled touch pins in turn to measure per START signal, i.e., only one measurement is conducted per START signal. As multiple START signals are generated, the Touch FSM will cycle through the enabled touch pins in sequential order.

When measuring for different touch sensors, the sleep interval between two measurements is \hyperref[fielddesc:PMUTOUCHWAITCYCLES]{PMU\_TOUCH\_WAIT\_CYCLES} + \hyperref[fielddesc:PMUTOUCHSLEEPCYCLES]{PMU\_TOUCH\_SLEEP\_CYCLES} with the clock cycle of LP\_DYN\_SLOW\_CLK.

\begin{tiplisting}
    \vspace{-2em}
    \begin{itemize}
        \item After the touch sensor receives the START signal from the touch timer, it waits a short period before starting the measurement to ensure the stability of the sampling signal. The wait time is configured by \hyperref[fielddesc:LPANATOUCHXPDWAIT]{LP\_ANA\_TOUCH\_XPD\_WAIT}. The timing clock is AON\_FAST\_CLK.
        \item If the touch timer is used to generate the START signal, the scanning can be conducted without software intervention.
    \end{itemize}
\end{tiplisting}

\subsubsection{Frequency Hopping}
\label{sec:touch-freq-hop}

When frequency hopping is enabled, the Touch FSM can measure the touch sensor at different frequency modes. Up to three frequency modes are supported, each corresponding to a set of configuration parameters. When a START signal is generated, the measurement will be carried out for each frequency mode sequentially, and the DONE signal will be generated when the measurements are completed for all frequency modes. Frequency hopping can be enabled together with the scan mode. In such case, users can monitor the touch condition of multiple touch sensors under different frequency modes by generating multiple START signals. The scanning process is shown in Figure \ref{fig:touch-freq-hop}.

\begin{figure}[H]
    \centering
    \includegraphics[width=\textwidth]{\modulefiles/figures/touch_freq_scan.pdf}
    \caption{Touch Sensor Frequency Hopping Sequence Timing}
    \label{fig:touch-freq-hop}
\end{figure}

Set \hyperref[fielddesc:LPANAFREQSCANEN]{LP\_ANA\_FREQ\_SCAN\_EN} to 1 to enable frequency hopping, configure \hyperref[fielddesc:LPANAFREQSCANCNTLIMIT]{LP\_ANA\_FREQ\_SCAN\_CNT\_LIMIT} to select the maximum number of frequency mode supported, and configure the following analog parameter registers to adjust the sampling performance of the touch sensor at different frequency modes:

\begin{itemize}
    \item \hyperref[fielddesc:LPANATOUCHFREQ0DBIAS]{LP\_ANA\_TOUCH\_FREQ\regindex{n}\_DBIAS} (\regindex{n} = 0-2)
    \item \hyperref[fielddesc:LPANATOUCHFREQ0DRVHS]{LP\_ANA\_TOUCH\_FREQ\regindex{n}\_DRV\_HS} (\regindex{n} = 0-2)
    \item \hyperref[fielddesc:LPANATOUCHFREQ0DRVLS]{LP\_ANA\_TOUCH\_FREQ\regindex{n}\_DRV\_LS} (\regindex{n} = 0-2)
    \item \hyperref[fielddesc:LPANATOUCHFREQ0DRESLPF]{LP\_ANA\_TOUCH\_FREQ\regindex{n}\_DRES\_LPF} (\regindex{n} = 0-2)
    \item \hyperref[fielddesc:LPANATOUCHFREQ0DCAPLPF]{LP\_ANA\_TOUCH\_FREQ\regindex{n}\_DCAP\_LPF} (\regindex{n} = 0-2)
\end{itemize}

Upon each START signal, the Touch FSM measures consecutively for frequency mode 0, 1, and 2. A START signal triggers a maximum of three measurement operations. When the measurements are completed for all frequency modes, the START signal is de-asserted until the next measurement. Therefore, the Touch FSM performs measurements on the enabled touch sensors in sequence upon multiple START signals. \tagged{ESP32-P4-latest}{A touch on the same sensor is recognized only when the number of hit frequency points is greater than or equal to \hyperref[fielddesc:LPANAFREQSCANCNTRISE]{LP\_ANA\_FREQ\_SCAN\_CNT\_RISE}.}

\begin{tiplisting}
    \vspace{-2em}
    \begin{itemize}
        \item When the same touch sensor is measured at different frequency modes, the wait time between measurements can be configured by \hyperref[fielddesc:LPANATOUCHXPDWAIT]{LP\_ANA\_TOUCH\_XPD\_WAIT} (count clock: AON\_FAST\_CLK).
        \item After receiving the START signal from the touch timer, the touch sensor waits for a short period of time before starting measuring to ensure the stability of the sampling signal. The wait time can be configured by \hyperref[fielddesc:LPANATOUCHXPDWAIT]{LP\_ANA\_TOUCH\_XPD\_\\WAIT} (count clock: AON\_FAST\_CLK).
        \item If a touch timer is used to generate the START signal, the scanning can be conducted without software intervention.
    \end{itemize}
\end{tiplisting}

\subsection{Touch Detection}

\subsubsection{Sampled Value}
\label{para:sampled-value}

As introduced in Section \ref{para:measurement-process}, the sampled value of a measurement can be read from \hyperref[fielddesc:RTCTOUCHPAD0DATA]{RTC\_TOUCH\_PAD\regindex{n}\_DATA}. However, the specific type of sampled value stored in \hyperref[fielddesc:RTCTOUCHPAD0DATA]{RTC\_TOUCH\_PAD\regindex{n}\_DATA} is determined together by \hyperref[fielddesc:LPANATOUCHDATASEL]{LP\_ANA\_TOUCH\_DATA\_SEL} and \hyperref[fielddesc:LPANATOUCHFREQSEL]{LP\_ANA\_TOUCH\_FREQ\_SEL}. The following three sampled value types are supported:

\begin{enumerate}
    \item \textbf{touch\_raw\_data}: When the measurement is completed, the touch counter's value indicates the time taken to charge and discharge the touch pin a fixed number of times in LP\_DYN\_SLOW\_CLK cycles. After the IIR filter is turned on, the data is the filtered data.

    \item \textbf{touch\_smooth\_data}: A series of moving averages generated by the touch\_raw\_data values sampled from a single touch sensor filtred by an IIR filter.

    touch\_smooth\_data is less prone to noise spikes or outlier samples, thus often used for hardware touch detection. If frequency hopping is enabled, each touch pin will have a separate set of smooth\_touch\_data at different frequency modes and do not affect each other. The IIR filter that generates touch\_smooth\_data is configured via \hyperref[fielddesc:LPANATOUCHSMOOTHLVL]{LP\_ANA\_TOUCH\_SMOOTH\_LVL}. See Table \ref{tab:smooth-cal} for more information.

    \begin{longtable}[c]{|p{5.5 cm}|p{1.5cm}|p{8.5cm}|}
        \caption{touch\_smooth\_data Calculation}
        \label{tab:smooth-cal}\\
        \hline
        \rowcolor{lightgray}
        \textbf{\hyperref[fielddesc:LPANATOUCHSMOOTHLVL]{LP\_ANA\_TOUCH\_SMOOTH\_LVL}} & \textbf{Type} & \textbf{Formula}\\ 
        \hline
        \endhead
        \endfoot
        \endlastfoot
        0  & --- & touch\_raw\_data  \\\hline
        1  & IIR 1/2 & 1/2 touch\_raw\_data + 1/2 touch\_smooth\_data \\\hline
        2  & IIR 1/4 & 1/4 touch\_raw\_data + 3/4 touch\_smooth\_data \\\hline
        3  & IIR 1/8 & 1/8 touch\_raw\_data + 7/8 touch\_smooth\_data \\\hline
    \end{longtable}

    \item \textbf{benchmark}: Also generated from the touch\_raw\_data values by IIR filtering. However, the moving average window for benchmark is much wider.
    
    The benchmark value is usually used to represent the stable reading of a touch pin without the effect from a touch action. The touch detection is based on the benchmark value. Each touch sensor has a separate set of benchmark values at different frequency modes with a default value of 0xFFFF. The IIR filter that generates benchmark is configured via \hyperref[fielddesc:LPANATOUCHFILTERMODE]{LP\_ ANA\_TOUCH\_FILTER\_MODE}. Refer to Table \ref{tab:benchmark-cal} for more information.

    \iftagged{ESP32-P4-latest}{
    In addition to the IIR filtering, users can also update the benchmark value directly by software. The software update procedure is as follows:
    \begin{itemize}
      \item Configure the target frequency point to be updated through \hyperref[fielddesc:LPANATOUCHUPDATEBENCHMARKFREQSEL]{LP\_ANA\_TOUCH\_UPDATE\_BENCHMARK\_FREQ\_SEL}.
      \item Configure the target channel to be updated through \hyperref[fielddesc:LPANATOUCHUPDATEBENCHMARKPADSEL]{LP\_ANA\_TOUCH\_UPDATE\_BENCHMARK\_PAD\_SEL}.
      \item Configure the benchmark data value through \hyperref[fielddesc:LPANATOUCHBENCHMARKSW]{LP\_ANA\_TOUCH\_BENCHMARK\_SW}.
      \item Configure \hyperref[fielddesc:LPANATOUCHUPDATEBENCHMARKSW]{LP\_ANA\_TOUCH\_UPDATE\_BENCHMARK\_SW} to trigger the benchmark data update by software.
    \end{itemize}
    
    This method is generally used for software debugging. When the benchmark data of the target frequency point and channel has been updated, the interrupt \hyperref[int:TOUCHBENCHMARKUPDATEINT]{TOUCH\_BENCHMARK\_UPDATE\_INT} will be generated.
    }
    {In addition to the IIR filtering, users can also update the benchmark value directly by software. Configure \hyperref[fielddesc:LPANATOUCHBENCHMARKSW]{LP\_ANA\_TOUCH\_BENCHMARK\_SW} and \hyperref[fielddesc:LPANATOUCHUPDATEBENCHMARKSW]{LP\_ANA\_TOUCH\_UPDATE\_BENCHMARK\_SW} sequentially to update the benchmark value. This is generally used for software debugging.}

\end{enumerate}

\begin{longtable}[c]{|p{5.5 cm}|p{1.5cm}|p{8.5cm}|}
    \caption{benchmark Algorithm}\label{tab:benchmark-cal}\\
    \hline
    \rowcolor{lightgray}
    \textbf{\hyperref[fielddesc:LPANATOUCHFILTERMODE]{LP\_ANA\_TOUCH\_FILTER\_MODE}} & \textbf{Type} & \textbf{Formula}\\ 
    \hline
    \endhead
    \endfoot
    \endlastfoot
    0  &  IIR 1/4 & 1/4 touch\_raw\_data + 3/4 benchmark \\\hline
    1  &  IIR 1/8 & 1/8 touch\_raw\_data + 7/8 benchmark \\\hline
    2  &  IIR 1/16 & 1/16 touch\_raw\_data + 15/16 benchmark \\\hline
    3  &  IIR 1/32 & 1/32 touch\_raw\_data + 31/32 benchmark \\\hline
    4  &  IIR 1/64 & 1/64 touch\_raw\_data + 63/64 benchmark \\\hline
    5  &  IIR 1/128 & 1/128 touch\_raw\_data + 127/128 benchmark \\\hline
    6  &  IIR 1/256 & 1/256 touch\_raw\_data + 255/256 benchmark \\\hline
    7  &  JITTER & touch\_raw\_data +/$–$ \hyperref[fielddesc:LPANATOUCHJITTERSTEP]{LP\_ANA\_TOUCH\_JITTER\_STEP} \\\hline
\end{longtable}

\subsubsection{Hardware Touch Detection}

In hardware touch detection, the touch sensor can detect finger touch or release and trigger an interrupt. To use hardware touch detection, the following key parameters need to be defined:

\begin{itemize}
    \item \textbf{finger\_threshold}: The threshold value used to determine a touch or a touch interruption. finger\_threshold is configured via \hyperref[fielddesc:LPANATOUCHPAD0TH0]{LP\_ANA\_TOUCH\_PAD\regindex{x}\_TH\regindex{y}}, where \regindex{x} (1-14) represents the touch pin index and \regindex{y} (0-2) represents the frequency mode configured in hopping frequency.

    \item \textbf{noise\_threshold}: Indicates the maximum threshold of the upward and downward fluctuation of touch\_smooth\_data. noise\_threshold is configured via \hyperref[fielddesc:LPANATOUCHNOISETHRES]{LP\_ANA\_TOUCH\_NOISE\_THRES}.
    
    \item \textbf{hysteresis}: Used to avoid touch misjudgments caused by fluctuations of touch\_smooth\_data around touch thresholds. The hysteresis value is configured via \hyperref[fielddesc:LPANATOUCHHYSTERESIS]{LP\_ANA\_TOUCH\_HHYSTERESIS}.

\end{itemize}

The finger\_threshold and noise\_threshold values are not absolute thresholds, but offsets from the benchmark value, while the hysteresis value is an offset from finger\_threshold. This prevents from misdiagnosing the gradual or slow changes of touch\_raw\_data caused by environmental factors such as temperature, power, or noise as a touch action.

After a series of calculations and comparisons based on the above parameters, the touch sensor will output the touch results detected on the touch pins via interrupts. Each touch sensor corresponds to a control register \hyperref[fielddesc:LPANATOUCHOUTEN]{LP\_ANA\_TOUCH\_OUTEN} to configure whether to output the detection result.

\subsection{Proximity Mode}

When an object such as a finger is close to, but not touching, a touch pin, a small change in capacitance occurs on the touch pin, which is much smaller than the change caused by a physical touch. Proximity mode allows for the detection of these small changes.

\begin{figure}[H]
    \centering
    \includegraphics[width=0.17\textwidth]{\modulefiles/figures/touch_approach.eps}
    \caption{Sensing Area}
    \label{fig:touch-approach}
\end{figure}

\begin{itemize}
    \item The maximum detection distance ``d" is 16 cm and the sensing area ``S" is 20 cm$^{2}$, where ``d" is positively correlated with ``S" as shown in Figure \ref{fig:touch-approach}.
    \item Up to three touch pins can be configured to operate simultaneously in proximity mode.
    \item Frequency hopping is also supported in proximity mode.
\end{itemize}

In proximity mode, a touch sensor samples for a fixed number of times and accumulates the sample values. If the final cumulative value exceeds the configured threshold, a proximity object is considered to be detected and an interrupt will be triggered. Note that since the sample values are accumulated, the values introduced in Section \ref{para:sampled-value}, i.e., touch\_raw\_data, touch\_smooth\_data, and benchmark, will not be generated in proximity mode.

To operate in proximity mode:

\begin{itemize}
    \item Configure a touch sensor to operate in proximity mode by setting \hyperref[fielddesc:LPANATOUCHAPPROACHPAD0]{LP\_ANA\_TOUCH\_APPROACH\_PAD0}, \hyperref[fielddesc:LPANATOUCHAPPROACHPAD1]{LP\_ANA\_TOUCH\_APPROACH\_PAD1}, or \hyperref[fielddesc:LPANATOUCHAPPROACHPAD2]{LP\_ANA\_TOUCH\_APPROACH\_PAD2}.
    \item Configure \hyperref[fielddesc:LPANATOUCHAPPROACHLIMIT]{LP\_ANA\_TOUCH\_APPROACH\_LIMIT} to adjust the number of samplings to generate the cumulative value. The touch sensor contains an internal sample counter to record the number of samplings.
    \item Set the threshold value via \hyperref[fielddesc:LPANATOUCHPAD0TH0]{LP\_ANA\_TOUCH\_PAD\regindex{x}\_TH\regindex{y}}, where \regindex{x} (1-14) represents the touch pin index and \regindex{y} (0-2) represents the frequency mode configured in hopping frequency.
    \item When the sample counter reaches the number of samplings configured in \hyperref[fielddesc:LPANATOUCHAPPROACHMEASNUM0]{LP\_ANA\_TOUCH\_APPROACH\_MEAS\_NUM\regindex{n}}:
    \begin{itemize}
        \item If the cumulative value exceeds the threshold value, an interrupt will be triggered.
        \item The sample counter and the cumulative value are reset to 0. The touch sensor restarts accumulating the sample values.
    \end{itemize}
\end{itemize}

\subsection{Sleep Mode}

To reduce power consumption, the touch sensor also supports sleep mode. In sleep mode, users can designate one touch sensor to operate actively (hereinafter referred to as the sleeping touch sensor) to monitor whether a touch occurs in real time, while the remaining 13 touch sensors are all powered down in sleep mode. Once the sleeping touch sensor detects a touch or a touch release, it will generate an interrupt to wake up the rest of the touch sensors. The sleeping touch sensor also supports frequency hopping and proximity mode. There is a separate set of threshold parameters for the sleeping touch sensor, which can be configured as follows:

\begin{itemize}
    \item Configure \hyperref[fielddesc:LPANATOUCHSLPPAD]{LP\_ANA\_TOUCH\_SLP\_PAD} to select a particular touch sensor as the sleeping touch sensor.
    \item Configure \hyperref[fielddesc:LPANATOUCHSLPTH0]{LP\_ANA\_TOUCH\_SLP\_TH0}, \hyperref[fielddesc:LPANATOUCHSLPTH0]{LP\_ANA\_TOUCH\_SLP\_TH1}, and \hyperref[fielddesc:LPANATOUCHSLPTH0]{LP\_ANA\_TOUCH\_SLP\_TH2} to set the touch threshold for the sleeping touch sensor at different frequency modes.
    \item Configure \hyperref[fielddesc:LPANATOUCHSLPAPPROACHEN]{LP\_ANA\_TOUCH\_SLP\_APPROACH\_EN} to enable the proximity mode for the sleeping touch sensor.
    \item Configure \hyperref[fielddesc:LPANATOUCHSLPCHANNELCLR]{LP\_ANA\_TOUCH\_SLP\_CHANNEL\_CLR} to clear the intermediate data for the sleeping touch sensor, such as benchmark, touch\_raw\_data, and the number of sampling.
\end{itemize}

\begin{tiplisting}
The configuration for the noise\_threshold, hysteresis, touch\_smooth\_data, and benchmark values are the same for the sleeping touch sensor in sleep mode as in the normal mode.
\end{tiplisting}

\subsection{Moisture Tolerance}

If there are water droplets on the touch pads and the droplets are large enough to physically bridge two more adjacent touch pads, the adjacent touch pads may be electrically coupled. Coupled touch pads can cause false touch detection due to capacitance changes caused by the coupling. If a water droplet is connected across the touch pad and ground, a large parasitic capacitance will be introduced, also leading to false detection. The moisture tolerance feature can mitigate the impact of water droplets.

To configure moisture tolerance:

\begin{itemize}
    \item Configure \hyperref[fielddesc:LPANATOUCHSHIELDPADEN]{LP\_ANA\_TOUCH\_SHIELD\_PAD\_EN} to enable moisture tolerance.
    \item Configure \hyperref[fielddesc:LPANATOUCHBUFSEL]{LP\_ANA\_TOUCH\_BUFSEL} to select a touch pin to be used for moisture tolerance.
\end{itemize}

\begin{tiplisting}
The moisture tolerance feature only works at frequencies below 10 MHz. Therefore, when the frequency is higher than 10 MHz, please disable the moisture tolerance feature.
\end{tiplisting}

\subsection{Water Rejection}
\label{sec:tchsen-water-reject}

If the sensor array becomes wet, i.e., the majority of the sensor pads are soaked by water, most (if not all) of the touch pads will become unusable due to the possible false touch detections. \chipname{} supports the water rejection feature to shut down the sensor array if it is detected to be wet.

Configure \hyperref[fielddesc:LPANATOUCHOUTRING]{LP\_ANA\_TOUCH\_OUT\_RING} to select one of the touch pins to be used for the water rejection feature.


%%%%%%%%%%%%%%%%%%%%%
%%% Interrupts %%%
%%%%%%%%%%%%%%%%%%%%%

\section{Interrupts}

\chipname{}'s touch sensor can generate the LP\_TOUCH\_INTR signal that will be sent to the \hyperref[mod:intmtrx]{Interrupt Matrix}.

There are several internal interrupt sources from touch sensor that can generate the LP\_TOUCH\_INTR signal. The interrupt sources from touch sensor are listed with their trigger conditions and the resulted interrupt signal in Table \ref{tab:tchsen-int-sources}.

\newpage
\begin{longtable}{ | L{6.2cm} | L{6.5cm} | L{3cm} | }
\caption{Touch Sensor's Internal Interrupt Sources} \label{tab:tchsen-int-sources}
\\ \hline
\rowcolor{lightgray}
\textbf{Internal Interrupt Source} & \textbf{Trigger Condition}  &  \textbf{Interrupt Signal}\\ \hline
\label{int:TOUCHAPPROACHLOOPDONEINT}TOUCH\_APPROACH\_LOOP\_DONE\_INT  &  Completion of proximity mode cumulative sampling &  \multirow{6}*{LP\_TOUCH\_INTR}  \\ \cline{1-2}
\label{int:TOUCHTIMEOUTINT}TOUCH\_TIMEOUT\_INT &  Timeout for the TOUCH\_OUT signal sampling & \\ \cline{1-2}
\label{int:TOUCHINACTIVEINT}TOUCH\_INACTIVE\_INT &  Touch release detected & \\ \cline{1-2}
\label{int:TOUCHACTIVEINT}TOUCH\_ACTIVE\_INT &  Touch detected &  \\ \cline{1-2}
\label{int:TOUCHDONEINT}TOUCH\_DONE\_INT & Completion of sampling an individual touch pin at a certain frequency mode for comparison &\\ \cline{1-2}
\iftagged{ESP32-P4-latest}{
\label{int:TOUCHSCANDONEINT}TOUCH\_SCAN\_DONE\_INT &  Completion of sampling selected touch pins for comparison at all frequency modes & \\ \cline{1-2}
\label{int:TOUCHBENCHMARKUPDATEINT}TOUCH\_BENCHMARK\_UPDATE\_INT &  Completion of software benchmark data update & \\ \hline}
{\label{int:TOUCHSCANDONEINT}TOUCH\_SCAN\_DONE\_INT &  Completion of sampling selected touch pins for comparison at all frequency modes & \\ \hline}
\end{longtable}

\begin{tiplisting}
For definitions of \textit{interrupt}, \textit{interrupt signal}, \textit{interrupt source}, and their correlations, please refer to Chapter \ref{mod:intmtrx} \textit{\nameref{mod:intmtrx}} > Section \ref{sec:intmtx-int-terminology} \textit{\nameref{sec:intmtx-int-terminology}}.
\end{tiplisting}

Each interrupt source can be configured by a common set of registers that are described in Section \hyperref[interrupt-config-registers]{\textit{Interrupt Configuration Registers}}. The specific registers can be found in Section \ref{sec:tchsen-reg-summ} \textit{\nameref{sec:tchsen-reg-summ}}.


% >>> If your module has no registers, delete sample text
%       for Sections Register Summary and Registers below <<<

%%%%%%%%%%%%%%%%%%%%%%%%
%%% Register Summary %%%
%%%%%%%%%%%%%%%%%%%%%%%%

% >>> Update the sample text below and insert
%     the info on register summary into the subfile <<<

\newpage
\begin{landscape}
\hypertarget{tchsen-reg-summ}{}
\section{Register Summary}
\label{sec:tchsen-reg-summ}

\subsection{Interrupt and Status Register Summary}
\label{sec:tchsen-status-reg-summ}

The addresses in this section are relative to Touch Sensor base address provided in Table \ref{tab:sysmem-base-address} in Chapter \ref{mod:sysmem} \textit{\nameref{mod:sysmem}}.

The abbreviations given in Column \textbf{Access} are explained in Section \hyperref[glossary-access-types]{\textit{Access Types for Registers}}.

\subfile{\modulefiles/61-TCHSEN-reg-summ__EN}

\subsection{Configuration Register Summary}
\label{sec:tchsen-config-reg-summ}

The addresses in this section are relative to LP Analog Peripheral base address provided in Table \ref{tab:sysmem-base-address} in Chapter \ref{mod:sysmem} \textit{\nameref{mod:sysmem}}.

The abbreviations given in Column \textbf{Access} are explained in Section \hyperref[glossary-access-types]{\textit{Access Types for Registers}}.

\subfile{\modulefiles/61-TCHSEN-LP-reg-summ__EN}
\end{landscape}


%%%%%%%%%%%%%%%%%
%%% Registers %%%
%%%%%%%%%%%%%%%%%

% >>> Update the sample text below and insert
%      the info on registers into the subfile <<<

\newpage
\section{Registers}

\subsection{Interrupt and Status Registers}
\label{sec:tchsen-status-reg}

The addresses in this section are relative to Touch Sensor base address provided in Table \ref{tab:sysmem-base-address} in Chapter \ref{mod:sysmem} \textit{\nameref{mod:sysmem}}.

For how to program reserved fields, please refer to Section \hyperref[programming-reserved-field]{\textit{Programming Reserved Register Field}}.

\subfile{\modulefiles/61-TCHSEN-reg__EN}

\subsection{Configuration Registers}
\label{sec:tchsen-config-reg}

The addresses in this section are relative to LP Analog Peripheral base address provided in Table \ref{tab:sysmem-base-address} in Chapter \ref{mod:sysmem} \textit{\nameref{mod:sysmem}}.

For how to program reserved fields, please refer to Section \hyperref[programming-reserved-field]{\textit{Programming Reserved Register Field}}.

\subfile{\modulefiles/61-TCHSEN-LP-reg__EN}

\end{document}
